\documentclass[11pt]{article}
\usepackage[a4paper,margin=21mm]{geometry}
\usepackage{fontspec}
\setmainfont{Latin Modern Roman}
\newfontfamily\CJKfont{Noto Serif CJK SC}
\XeTeXlinebreaklocale "zh"
\XeTeXlinebreakskip=0pt plus 1pt
\usepackage{amsmath,amssymb,amsthm,amscd,graphicx,color}
\usepackage{xurl}
\usepackage[unicode,hidelinks]{hyperref}
\allowdisplaybreaks
\setlength\emergencystretch{3em}
\numberwithin{equation}{section}
\usepackage{mathrsfs,mathtools,etoolbox}
% Independently written finite standard-notation bindings, 1170 only.
% No gasymbols implementation is copied or executed.
\newcommand{\N}{\mathbb{N}}
\newcommand{\Z}{\mathbb{Z}}
\newcommand{\dist}{\mathrm{d}}
\newcommand{\dif}{\,\mathrm{d}}
\newcommand{\dd}{\mathrm{d}}
\newcommand{\D}{\mathrm{D}\kern.01cm}
\newcommand{\CS}{\mathscr{C}}
\newcommand{\BV}{\mathrm{BV}}
\DeclareMathOperator{\Int}{Int}
\newcommand{\ncapa}{\mathfrak{c}}
\DeclareMathOperator{\iso}{iso}
\let\div\relax
\DeclareMathOperator{\div}{div}
\DeclareMathOperator{\loc}{loc}
\DeclareMathOperator{\ee}{e}
\DeclareMathOperator{\Crit}{Crit}
\let\H\relax
\DeclareMathOperator{\H}{H}
\DeclareMathOperator{\h}{h}
\makeatletter
\newcommand{\kst}{\@ifnextchar[{\PackageError{sigma1170}{Unsupported optional kst argument}{Only observed default C is supported}}{\mathrm{C}}}
\makeatother
\DeclarePairedDelimiter{\sigmaAbs}{\lvert}{\rvert}
\newcommand{\abs}[1]{\ifstrempty{#1}{\sigmaAbs{\kern1pt\cdot\kern1pt}}{\mathchoice{\sigmaAbs*{#1}}{\sigmaAbs{#1}}{\sigmaAbs{#1}}{\sigmaAbs{#1}}}}
\newcommand{\set}[1]{\ifstrempty{#1}{\PackageError{sigma1170}{Unsupported empty set argument}{All original selected arguments are nonempty}}{\begingroup\def\st{\,\middle\vert\,}\left\lbrace#1\right\rbrace\endgroup}}

\usepackage{tikz}
\usepackage{mathrsfs}
\usepackage{mathtools,leftidx}
\usepackage[shortcuts]{extdash}

\numberwithin{equation}{section}

\newtheorem{Theorem}{Theorem}[section]
\newtheorem{Lemma}[Theorem]{Lemma}
\newtheorem{Fact}[Theorem]{Fact}
\newtheorem{Corollary}[Theorem]{Corollary}
\newtheorem{Proposition}[Theorem]{Proposition}
\newtheorem{Claim}[Theorem]{Claim}

 { \theoremstyle{definition}
\newtheorem{Definition}[Theorem]{Definition}
\newtheorem{Remark}[Theorem]{Remark} }


\makeatletter
\newenvironment{claim}[1][]{\par
 %
 \normalfont \topsep7\p@\@plus6\p@\relax
 \trivlist
%<amsbook|amsproc> \itemindent\normalparindent
 \item[\indent\hskip\labelsep
%<amsbook|amsproc> \scshape
%<amsart|amsthm> \itshape
 \ifx\relax#1\relax \textbf{Claim}\@addpunct{.}\else#1\@addpunct{.}\fi]\ignorespaces \itshape
}{%
 \normalfont\endtrivlist\@endpefalse
}

\newcounter{case}
\@addtoreset{case}{proof}

\newenvironment{case}[1][]{\refstepcounter{case}\par
 %
 \normalfont \topsep7\p@\@plus6\p@\relax
 \trivlist
 \medskip
%<amsbook|amsproc> \itemindent\normalparindent
 \item[\itshape\hskip\labelsep
%<amsbook|amsproc> \scshape
%<amsart|amsthm> \itshape
 Case \thecase\@addpunct{.} #1\@addpunct{.}]\ignorespaces
}{%
 \endtrivlist\@endpefalse
 \medskip
}
\makeatother

%\newcommand{\crefitem}[2]{\cref{#1}\cref{#2}}


\makeatletter
\newcommand{\leqnmode}{\tagsleft@true}
\makeatother
\renewcommand{\ncapa}{\mathfrak{c}}
\makeatletter
\newcommand{\svm}{\gamma}
\makeatother
\newcommand{\eum}{\delta}
\newcommand{\ma}{\mathfrak{m}}
\newcommand{\A}{\mathcal{A}}
\newcommand{\tp}[1][p]{{\text{\tiny$(#1)$}}}


\newcommand{\lop}{\mathscr{L}}

\makeatletter\newlabel{eq:IMCF}{{1.2}{}{Original source locator}{}{}}\makeatother
\makeatletter\newlabel{eq:isoperimetric_mass_intro}{{1.1}{}{Original source locator}{}{}}\makeatother
\begin{document}
\begin{center}\Large Nonsharp p-isocapacitary Penrose inequality with a nonempty connected minimal boundary, vanishing relative H2 and the stated level-gradient decay\end{center}
\noindent\textbf{Stable ID:} 1170\quad\textbf{Author:} Luca Benatti; Mattia Fogagnolo; Lorenzo Mazzieri\par
\noindent\textbf{Work:} Nonlinear Isocapacitary Concepts of Mass in3-Manifolds with Nonnegative Scalar Curvature\par
\noindent\textbf{Source:} \url{https://sigma-journal.com/2023/091/}\par
\noindent\textbf{Version:} Official SIGMA 19 (2023) article 091, DOI 10.3842/SIGMA.2023.091; journal PDF and native TeX acquired 2026-09-30, exact bytes pinned by SHA256\par
\noindent\textbf{Rights:} CC BY 4.0. Authors retain copyright. \url{https://creativecommons.org/licenses/by/4.0/}. Reuse and typesetting adaptation; original language and mathematical content preserved.\par
\noindent\textbf{Difficulty (prior selection preserved):} {\CJKfont H3: The proof develops weak monotonicity of two nonlinear p-Hawking quantities and compares their common asymptotic limit to the volume-capacity mass. It regularizes the p-capacitary equation, passes second-derivative terms distributionally through critical sets, and controls the residual terms by weak Sobolev convergence and dominated convergence. The comparison uses a two-sign asymptotic subsequence argument and combines both monotonic quantities to obtain the precise nonsharp coefficient (5-p)/2.}\par
\medskip\noindent\textbf{Editorial boundary note (not author text):} Selected author claim is precisely inequality(1.5) of Theorem1.2 for a nonempty smooth compact connected minimal boundary, with all smooth/connected/complete/noncompact/one-ended, strong p-nonparabolicity,1<p<3, scalar-curvature, relative-H2 and exact dagger level-gradient assumptions retained. Full native Theorem1.2 is preserved as source context, but the empty-boundary positivity/Euclidean rigidity sentences are outside the counted obligation, and their proof is unselected. The complete current two-mass comparison, both signs of the limiting subsequence, both AppendixA full monotonicity/distributional regularization proofs and every supporting definition are included. Existing p-harmonic regularity, precise regularized differential inequality and standard generalized l-Hospital criterion remain author references. One independently announced first-main nonsharp Penrose outcome; no extra count for supporting monotonicity/comparison. A bounded independently written standard AMS/mathtools notation wrapper binds exactly the20 actually used names; native author expressions are unchanged. The unidentified gasymbols.sty is never copied, executed or exported, and the original ZIP containing it stays private and hash-bound. Exactly two empty absolute-value arguments retain dot-in-bars notation, all42 set arguments remain nonempty with grouped separators, and kst rejects optional arguments.\par
\noindent\textbf{External original reference eq:IMCF:} Original standard weak inverse-mean-curvature comparison flow, source146–154; mentioned only as a contrasting prior flow in selected introductory context, not used in the counted nonlinear proof.\par
\noindent\textbf{External original reference eq:isoperimetric\_mass\_intro:} Original prior isoperimetric-mass definition, source133–139, mentioned only as the comparator motivating the fully retained p-isocapacitary definition; not a missing step in the counted nonlinear mass proof.\par
\medskip\hrule\medskip\noindent\textbf{Author text begins below.}\par
\setcounter{section}{1}\setcounter{equation}{2}
% Original source lines 156--212
In the present paper, we are going to develop a similar theory, in the case where the weak IMCF~\eqref{eq:IMCF} is replaced by the level set flow of weak solutions $w_p\in \CS^{1,\beta}_{\loc}(M \smallsetminus \Omega)$ to the boundary value problem
\begin{equation}
\label{eq:pb-intro}
 \begin{cases}
 \Delta_p w_p = \abs{\D w_p}^p &\text{on $M \smallsetminus \Int \Omega$,}\\
 w_p=0 & \text{on $\partial \Omega$,}\\
 w_p\to +\infty & \text{as $\dist(x,\partial \Omega) \to +\infty$.}
 \end{cases}
\end{equation}
The link between the above problem and the weak IMCF~\eqref{eq:IMCF} relies on the fact that $w_p\to w_1$ as $p\to 1^+$ locally uniformly on $M$~\cite{kotschwar_localgradientestimatesharmonic_2009,mari_flowlaplaceapproximationnew_2022, moser_inversemeancurvatureflow_2007,moser_inversemeancurvatureflow_2008}, provided some natural global requirements are met by the manifold $(M, g)$. On the other hand,
the solutions to problem~\eqref{eq:pb-intro} are directly connected to the notion of $p$-capacitary potential of a compact body $\Omega$. In fact, setting $w_p = -(p-1) \log u_p$ implies that $u_p$ is $p$-harmonic, that is $\Delta_p u_p = 0$. These relationships have been instrumental in demonstrating a series of geometric inequalities by means of monotonicity formulas, holding along the level sets of solutions to equation~\eqref{eq:pb-intro}. As $p \to 1^+$, these inequalities become increasingly close to the desired result.
This machinery, firstly introduced in the case $p=2$ for harmonic functions in~\cite{agostiniani_sharpgeometricinequalitiesclosed_2020, agostiniani_monotonicityformulaspotentialtheory_2020}, has proven to be powerful enough to produce an enhanced version of the Minkowski inequality~\cite{agostiniani_minkowskiinequalitiesnonlinearpotential_2022, fogagnolo_geometricaspectscapacitarypotentials_2019}, later extended to Riemannian manifolds with nonegative Ricci curvature~\cite{benatti_minkowskiinequalitycompleteriemannian_2022} as well as to the anisotropic setting~\cite{xia_anisotropiccapacityanisotropicminkowski_2022}.
 In~\cite{agostiniani_greenfunctionproofpositive_2021} and subsequently in~\cite{agostiniani_riemannianpenroseinequalitynonlinear_2022}, the authors used this approach on $3$-manifolds with nonnegative scalar curvature to prove the Riemannian Penrose inequality for a single black hole, based on the monotonic behaviour of a suitable $p$-harmonic version of the Hawking mass (see~\eqref{eq:p-Hawkingmass} below).
\emph{Throughout the manuscript, Riemannian manifolds are assumed to be smooth, connected, metrically complete, noncompact, with one single end.}


The main object of interest in the present paper is the following nonlinear potential theoretic version of Huisken's isoperimetric mass~\eqref{eq:isoperimetric_mass_intro}, that we call the {\em $p$-isocapacitary mass}. It involves the classical notion of $p$-capacity of a compact set $K \subset M$, that, in dimension $3$ and for $1 < p <3$, is given by
\begin{equation*}%\label{eq:p-capacity-intro}
 \ncapa_p(K) = \inf \set{\frac{1}{4\pi}\bigg(\frac{p-1}{3-p}\bigg)^{p-1}\int_{M\smallsetminus K} \abs{\D v}^p \dif \mu\st v \in \CS_c^\infty(M),\, v \geq 1\text{ on } K}.
\end{equation*}
When the boundary of $K$ is regular enough, such infimum is realized by the $p$-capacitary potential of $K$, i.e., the unique $p$-harmonic function $u_p$, equal to $1$ on $\partial K$ and vanishing at infinity.
\begin{Definition}[$p$-isocapacitary mass]%\label{def:iso-p-capacitary_mass}
Let $(M,g)$ be a Riemannian $3$-manifold with compact boundary, and let $1<p<3$. Given a closed bounded subset $\Omega\subset M$ containing $\partial M$ with $\CS^{1}$-boundary the quasi-local $p$-isocapacitary mass of $\Omega$ is defined as
\begin{equation*}%\label{eq:quasi-local_isopcapacitary_mass}
 \ma_{\iso}^{\tp}(\Omega)=\frac{1}{2p \pi \ncapa_p(\partial \Omega)^{\frac{2}{3-p}}}\left( \abs{\Omega} -\frac{4 \pi}{3} \ncapa_p (\partial \Omega)^\frac{3}{3-p}\right).
\end{equation*}
The $p$-isocapacitary mass $\ma^{\tp}_{\iso}$ of $(M,g)$ is defined as \begin{equation}\label{eq:isopcapacitary_mass}
 \ma_{\iso}^{\tp} = \sup_{(\Omega_j)_{j \in \N}} \limsup_{j \to +\infty} \ma_{\iso}^{\tp}(\partial \Omega_j),
\end{equation}
where the supremum is taken among all exhaustions $\set{\Omega_j}_{j\in \N}$.
\end{Definition}
As for the isoperimetric mass, the $p$-isocapacitary mass might be any number in $[-\infty,+\infty]$. The special and particularly relevant case of the $2$-isocapacitary mass has been recently introduced and studied by Jauregui~\cite{jauregui_admmasscapacityvolumedeficit_2020}.

A first natural and fundamental question about these newly introduced quantities is whether they are nonnegative, on the class of $3$-manifolds with nonnegative scalar curvature, where a~solution to~\eqref{eq:pb-intro} exists for any bounded $\Omega$ with regular boundary. The latter mentioned property will be referred to as the \textit{strong $p$-nonparabolicity} of the manifold $(M,g)$, a terminology that interpolates between the notion of strong nonparabolicity introduced by Ni~\cite{ni_meanvaluetheoremsmanifolds_2007} and the notion of strong $1$-nonparabolicity, which was employed in~\cite{benatti_isoperimetricriemannianpenroseinequality_2022} and recalled above.

Our first main result is a nonlinear potential-theoretic version of the Riemannian Penrose inequality, that, although not sharp, implies the positive mass theorem for the $p$-isocapacitary mass, with its associated rigidity statement. We prove its validity under the following asymptotic integral gradient estimate:
\bgroup\leqnmode
\begin{equation*}
\begin{minipage}[t]{.9585\textwidth}%
Given any $\Omega \subset M$ closed bounded subset with smooth and connected boundary homologous to $\partial M$, the function $w_p\in \CS^1_{\loc}\big(M \smallsetminus \overline{\Omega}\big)$ that solves~\eqref{eq:pb-intro} satisfies
\begin{gather*}
\int_{\partial \Omega_t} \abs{ \D w_p}^2 \dif \sigma = o\bigl(\ee^{t/(p-1)}\bigr)
\end{gather*}
as $t \to +\infty$ where $\Omega_t= \set{ w_p \leq t}$.
\end{minipage}\tag{$\dag$}\label{eq:energy_hp}
\end{equation*}
\egroup

\begin{Theorem}[$p$-isocapacitary Riemannian Penrose inequality]
\label{thm:pmt-intro}
 Let $(M,g)$ be a strongly $p$-nonparabolic Riemannian $3$-manifold, $1<p<3$, with nonnegative scalar curvature and with smooth, compact, connected, minimal, possibly empty boundary. Assume also that $(M,g)$ satisfies~\eqref{eq:energy_hp} and $H_2( M, \partial M; \Z)=\set{0}$. Then
 \begin{equation}
\label{eq:p-penroseintro-crociata}
\ncapa_p(\partial M)^{\frac{1}{3-p}}\leq \frac{5-p}{2}\ma^{\tp}_{\iso}.
\end{equation}
In particular, $\ma^{\tp}_{\iso}\geq 0$ and it vanishes if and only if $(M,g)$ is isometric to the flat $3$-dimensional Euclidean space.
\end{Theorem}

\setcounter{section}{1}
% Original source lines 259--293
\section{Preliminaries in nonlinear potential theory}
\label{sec:2}
We first remind that the operator $\Delta_p$ is defined by $\Delta_p f = \div \big(\abs{\nabla f}^{p-2}\nabla f\big)$, for $f \in \CS^2$ with nonvanishing gradient.
As far as basic principles and regularity for \emph{weakly} $p$-harmonic functions are concerned, we just refer the reader to
\cite{dibenedetto_alphalocalregularityweak_1983,evans_newprooflocal1_1982,heinonen_asuperharmonicfunctionssupersolutionsdegenerate_1988,ladyzenskaja_linearquasilinearellipticequations_1968,lewis_regularityderivativessolutionscertain_1983,lieberman_boundaryregularitysolutionsdegenerate_1988,serrin_localbehaviorsolutionsquasilinear_1964,tolksdorf_dirichletproblemquasilinearequations_1983,trudinger_harnacktypeinequalitiestheir_1967,uraltseva_degeneratequasilinearellipticsystems_1968,valtorta_laplaceoperatorriemannianmanifolds_2013}, and~\cite{holopainen_nonlinearpotentialtheoryquasiregular_1990,holopainen_volumegrowthgreenfunctions_1999} for the theory of $p$-capacitary potentials in exterior domains (see also~\cite[Chapter 1]{benatti_monotonicityformulasnonlinearpotential_2022} and reference therein), including existence issues. The function $w_p$ that solves~\eqref{eq:pb-intro} can be written as $w_p= -(p-1)\log u_p$, where $u_p \in \CS^{1,\beta}_{\loc}(M \smallsetminus \Int \Omega)$ is the solution to
\begin{equation}\label{eq:p-cap-potential}
 \begin{cases}
 \Delta^{\tp}_g u_p = 0 &\text{on $M \smallsetminus \Int \Omega$,}\\
 u_p=1 & \text{on $\partial \Omega$,}\\
 u_p\to 0 & \text{as $\dist(x,\partial \Omega)\to +\infty$,}
 \end{cases}
\end{equation}
the following definition of strongly $p$-nonparabolic Riemannian manifold is consistent with that of strong nonparabolicity~\cite{ni_meanvaluetheoremsmanifolds_2007} and with the limit case of strong $1$-nonparabolicity~\cite{benatti_isoperimetricriemannianpenroseinequality_2022}.

\begin{Definition}[strongly $p$-nonparabolic]
We say that a $3$-dimensional Riemannian manifold~$(M,g)$ with compact possibly empty boundary is \emph{strongly $p$-nonparabolic}, $1<p<3$, if there exists a solution to~\eqref{eq:pb-intro} for some $\Omega \subseteq M$ closed bounded with smooth boundary homologous to $\partial M$.
\end{Definition}

\begin{Remark}
By the maximum principle, in a strongly $p$-nonparabolic manifold every $\Omega$ with $\CS^{1}$-boundary homologous to $\partial M$ admits a solution to~\eqref{eq:pb-intro}.
\end{Remark}

This definition naturally comes with the notion of the $p$-capacity of a compact subset $K \subset M$, which we recall to be
\begin{equation*}%\label{eq:p-capacity}
 \ncapa_p(K) = \inf \set{\frac{1}{4\pi}\bigg(\frac{p-1}{3-p}\bigg)^{p-1}\int_{M\smallsetminus K} \abs{\D v}^p \dif \mu\st v \in \CS_c^\infty(M),\, v \geq 1\text{ on } K}.
\end{equation*}
When the boundary of $K$ is sufficiently regular, the above infimum is realized by the function~$u_p$ solving~\eqref{eq:p-cap-potential}. A fundamental property of the $p$-capacity is that it is monotone with respect to the standard inclusion of sets. More specifically, if we consider sublevel sets $\Omega_t=\set{w_p \leq t}$ of solutions to~\eqref{eq:pb-intro}, we have that their $p$-capacities grow exponentially with respect to the arrival time parameter $t$. This is
completely analogous to the exponential growth of the area along the IMCF. We recall this useful property in the following lemma, proved in~\cite[Lemma~3.8]{holopainen_nonlinearpotentialtheoryquasiregular_1990}.

\begin{Lemma}
Let $(M,g)$ be a $3$-dimensional Riemannian manifold with compact possibly empty boundary $\partial M$. Let $\Omega\subseteq M$ be a closed bounded subset homologous to $\partial M$ with $\CS^1$-boundary, $w_p$ the solution to~\eqref{eq:pb-intro} starting at $\Omega$ and $\Omega_t= \set{w_p \leq t}$, $1<p<3$. We have
\[
 \ncapa_p(\partial \Omega_t) = \ee^{t} \ncapa_p(\partial \Omega)= \frac{1}{4\pi}\int_{\partial \Omega_t} \bigg(\frac{ \abs{ \D w_p}}{3-p}\bigg)^{p-1}\dif \sigma.
\]
\end{Lemma}

\setcounter{Theorem}{10}\setcounter{subsection}{1}\setcounter{equation}{10}
% Original source lines 461--481
\subsection{Concepts of mass in nonlinear potential theory}
The classical Hawking mass $\ma_H$, that is
\begin{equation}\label{eq:hawking-mass}
 \ma_{H}(\partial \Omega)= \frac{ \abs{ \partial \Omega}^{\frac{1}{2}}}{16 \pi^{\frac{3}{2}}} \bigg( 4 \pi - \int_{\partial \Omega}\frac{\H^2}{4} \dif \sigma \bigg),
\end{equation}
for $\Omega\subset M$, monotonically increases along the level sets of the weak IMCF~\cite{huisken_inversemeancurvatureflow_2001}. Such a property is clearly not preserved in general when one replaces the weak IMCF with solutions $w_p$ to~\eqref{eq:pb-intro}. For this reason, we consider a different family of quasi-local masses. We will call \emph{$p$-Hawking mass} the quantity
\begin{equation}\label{eq:p-Hawkingmass}
 \ma_H^{\tp}(\partial \Omega) = \frac{\ncapa_p(\partial \Omega)^{\frac{1}{3-p}}}{8 \pi}\Bigg[4\pi + \int_{\partial \Omega}\frac{\abs{\D w_p}^2}{(3-p)^2}\dif \sigma - \int_{\partial \Omega}\frac{\abs{ \D w_p} }{(3-p)}\H \dif \sigma \Bigg]
\end{equation}
for $\partial \Omega \in \CS^{1}$ with weak second fundamental form in $L^2(\partial \Omega)$. This should be thought of as a $p$-version of the classical Hawking mass. In fact, it is immediately seen that the $p$-Hawking mass formally converges to the Hawking mass as $p \to 1^+$, having in mind that along the weak IMCF $w_1$ we have $\abs{\D w_1} = \H$ and that the $p$-capacity of an outward minimizing set recovers the perimeter in such limit~\cite[Theorem 1.2]{fogagnolo_minimisinghullscapacityisoperimetric_2022}.
Crucially, as the Hawking mass is monotone along the weak IMCF~\cite[Geroch Monotonicity Formula 5.8]{huisken_inversemeancurvatureflow_2001},
 so the function $t \mapsto \ma_H^{\tp}(\partial \Omega_t)$, for $\Omega_t= \set{w_p \leq t}$, is monotone nondecreasing, as proven in~\cite{agostiniani_riemannianpenroseinequalitynonlinear_2022} (see actually~Appendix~\ref{app:monotonicities} for the \emph{full} monotonicity result).
\begin{Theorem}\label{thm:AMMO_monotonicity}
 Let $(M,g)$ be a strongly $p$-nonparabolic Riemannian $3$-manifold, ${1<p<3}$, with nonnegative scalar curvature and with connected, compact, possibly empty boundary. Assume that $H_2(M, \partial M;\Z) = \set{0}$. Let $\Omega\subseteq M$ be a closed bounded subset with connected $\CS^1$-boundary homologous to $\partial M$ and with second fundamental form $\h \in L^2(\partial \Omega)$, $w_p$ the solution to~\eqref{eq:pb-intro} starting at $\Omega$ and $\Omega_t= \set{w_p \leq t}$. The function $t\mapsto \ma_H^{\tp}(\partial \Omega_t)$ defined in~\eqref{eq:p-Hawkingmass} admits a~monotone nondecreasing $\BV_{\loc}(0,+\infty)$ representative and
 \begin{gather}
 \frac{\dd}{\dd t} \ma^{\tp}_H (\partial \Omega_t)=\frac{\ncapa_p(\partial \Omega_t)^{\frac{1}{3-p}}}{(3-p) 8 \pi} \Bigg( 4\pi -\int_{\partial \Omega_t}\frac{\mathrm{R}^\top}{2} \dif \sigma +\int_{\partial \Omega_t} \frac{\vert \mathring{\h} \vert^2}{2} + \frac{ \mathrm{R}}{2} + \frac{ \abs{ \D^\top \abs{ \D w_p} }^2}{\abs{ \D w_p}^2}\dif \sigma \nonumber \\
 \hphantom{\frac{\dd}{\dd t} \ma^{\tp}_H (\partial \Omega_t)=\frac{\ncapa_p(\partial \Omega_t)^{\frac{1}{3-p}}}{(3-p) 8 \pi} \bigg(}{}
 + \int_{\partial \Omega_t}\frac{5-p}{p-1}\bigg(\frac{ \abs{\D w_p}}{3-p} - \frac{ \H}{2}\bigg)^2\dif \sigma\Bigg)\label{eq:AMMO_monotonicity}
 \end{gather}
 holds at every $t$ regular for $w_p$.
\end{Theorem}

\setcounter{section}{2}
% Original source lines 540--637
\section[p-isocapacitary Riemannian Penrose inequality]{$\boldsymbol{p}$-isocapacitary Riemannian Penrose inequality}
\label{sec:3}
In establishing the asymptotic comparison between the $p$-Hawking mass~\eqref{eq:p-Hawkingmass} and the $p$-iso\-capacitary mass, the quantity
\begin{equation}\label{eq:p-geroch_modified} \tilde{\ma}^{\tp}_H(\partial \Omega)= \frac{\ncapa_p(\partial \Omega)^{\frac{1}{3-p}}}{4\pi (3-p)}\Bigg( 4 \pi - \int_{\partial \Omega} \frac{\abs{ \D w_p}^2}{(3-p)^2} \dif \sigma\Bigg)
\end{equation}
will naturally appear. This quantity is closely related to the ones studied in~\cite{chan_monotonicitygreenfunctions_2022,munteanu_comparisontheorems3dmanifolds_2023} (see Theorem~\ref{thm:CCLT_monotonicity} and Remark~\ref{rmk:cclt_mon} for details). In the following lemma, we discuss the monotonicity properties of~\eqref{eq:p-geroch_modified} along the level sets of the solution $w_p$ to~\eqref{eq:pb-intro} and its relations with the $p$-Hawking mass~\eqref{eq:p-Hawkingmass}.

\begin{Lemma}\label{thm:inequality_modifiedgeroch_AMMO}
Let $(M,g)$ be a strongly $p$-nonparabolic Riemannian $3$-manifold, $1<p<3$, with nonnegative scalar curvature and with connected, compact, possibly empty boundary. Assume also that $(M,g)$ satisfies~\eqref{eq:energy_hp} and $H_2( M, \partial M; \Z)=\set{0}$. Let $\Omega\subseteq M$ be a closed bounded subset with connected $\CS^1$-boundary homologous to $\partial M$ and with second fundamental form $\h \in L^2(\partial \Omega)$, $w_p$ the solution to~\eqref{eq:pb-intro} starting at $\Omega$ and $\Omega_t= \set{w_p \leq t}$. Then the function $t \mapsto \tilde{\ma}^{\tp}_H(\partial \Omega_t)$ belongs to $W^{1,1}_{\loc}(0,+\infty)$ is monotone nondecreasing. Moreover, we have
\begin{equation}\label{eq:inequality_modifiedgeroch_AMMO}
 \ma^{\tp}_{H} ( \partial \Omega_t)\leq\tilde{\ma}^{\tp}_H(\partial \Omega_t)
\end{equation}
for every $t \in [0,+\infty)$, and
\begin{equation}\label{eq:limit_behaviour_phawking}
 \lim_{t \to+\infty} \ma^{\tp}_{H} ( \partial \Omega_t) = \lim_{t \to +\infty} \tilde{\ma}^{\tp}_H(\partial \Omega_t).
\end{equation}
\end{Lemma}

\begin{proof}
Denote
\begin{gather*}
 N(t) =\ncapa_p(\partial \Omega_t)^{-\frac{1}{p-1}}\Bigg( 4 \pi - \int_{\partial \Omega_t} \frac{\abs{ \D w_p}^2}{(3-p)^2} \dif \sigma\Bigg), \\
 D(t)= \ncapa_p(\partial \Omega_t)^{-\frac{2}{(3-p)(p-1)}},
\end{gather*}
we have that
\[
 \tilde{\ma}^{\tp}_H(\partial \Omega_t)=(\ncapa_p(\partial \Omega_t)^{\frac{1}{3-p}}\Bigg( 4 \pi - \int_{\partial \Omega_t} \frac{\abs{ \D w_p}^2}{(3-p)^2} \dif \sigma\Bigg) =\frac{N(t)}{D(t)}.
\]
Observe that, $N(t)$ is the quantity studied in~\cite{chan_monotonicitygreenfunctions_2022,munteanu_comparisontheorems3dmanifolds_2023}, while $1/D(t)$ is an exponentially growing term we multiplied it by. The function $N(t)$, and thus $\tilde{\ma}^{\tp}_H(\partial \Omega_t)=N(t)/D(t)$, belongs to~$W^{1,1}_{\loc}(0,+\infty)$ by Theorem~\ref{thm:CCLT_monotonicity}. Moreover,
\eqref{eq:weak_derivative} in Theorem~\ref{thm:CCLT_monotonicity} gives that $N'(t)/D'(t) = 4\pi (3-p)\ma_{H}^{\tp}(\partial \Omega_t)$ which is nondecreasing by Theorem~\ref{thm:AMMO_monotonicity}. Finally, $D(t)\to 0$ as $t \to +\infty$, while $N(t) \to 0$ as $t \to +\infty$ by the assumption \eqref{eq:energy_hp}. It is now a general fact that given two functions $f,g\in W^{1,1}_{\loc}(0,+\infty)$, with $g(t)\neq 0$ and $g'(t)\neq 0$, such that $f(t),g(t)\to 0$ as $t \to +\infty$ and with $f'(t)/g'(t)$ monotone nondecreasing, the function $f(t)/g(t)$ is monotone nondecreasing as well (see, e.g.,~\cite[Lemma 3.2]{zhu_comparisongeometryriccicurvature_1997}). Applying it with $f(t)=N(t)$ and $g(t)=D(t)$, we have that the function in $t \mapsto \tilde{\ma}^{\tp}_H(\partial \Omega_t)$ is nondecreasing. To prove~\eqref{eq:inequality_modifiedgeroch_AMMO}, it is enough to observe that
\[
 0\leq \frac{ \dd }{\dd t} \left( \frac{N(t)}{D(t)}\right) = \frac{N'(t)D(t) - N(t)D'(t)}{D(t)^2} = \frac{8 \pi}{(p-1)}\bigl(-\ma^{\tp}_H(\partial \Omega_t) + \tilde{\ma}^{\tp}_{H}(\partial \Omega_t)\bigr)
\]
for almost any $t$.
It then remains to prove~\eqref{eq:limit_behaviour_phawking}. On the other hand, by de L'H\^opital rule (see~\cite[Theorem A.1]{benatti_isoperimetricriemannianpenroseinequality_2022})
\[
 \lim_{t \to +\infty} 4\pi (3-p) \tilde{\ma}^{\tp}_H(\partial \Omega_t)= \lim_{t \to +\infty} \frac{N(t)}{D(t)} \leq \lim_{t \to +\infty} \frac{N'(t)}{D'(t)} = \lim_{t \to +\infty} 4 \pi (3-p) \ma_H^{\tp}(\partial \Omega_t).
\]
The reverse inequality easily follows by~\eqref{eq:inequality_modifiedgeroch_AMMO}. \end{proof}

The following result gives the $p$-capacitary counterpart of~\cite[Lemma 2.7]{benatti_isoperimetricriemannianpenroseinequality_2022}, that asymptotically controls the Hawking mass with the quasi-local isoperimetric mass of the evolving sets.

\begin{Lemma}[asymptotic comparison lemma]\label{thm:mass_control_pcap}
Let $(M,g)$ be a strongly $p$-nonparabolic Riemannian $3$-manifold, $1<p<3$, with nonnegative scalar curvature and with connected, compact, possibly empty boundary. Assume that $H_2( M, \partial M; \Z)=\set{0}$ and $(M,g)$ satisfies~\eqref{eq:energy_hp}. Let $\Omega\subseteq M$ be a closed bounded subset with connected $\CS^1$-boundary homologous to $\partial M$ and with second fundamental form $\h \in L^2(\partial \Omega)$, $w_p$ the solution to~\eqref{eq:pb-intro} starting at $\Omega$ and $\Omega_t= \set{w_p \leq t}$. Then
\begin{equation}\label{eq:mass_contro_pcap}
 \lim_{t \to +\infty} \ma_{H}^{\tp}(\partial \Omega_t)= \lim_{t \to +\infty} \tilde{\ma}_{H}^{\tp}(\partial \Omega_t) \leq\liminf_{t \to +\infty} \ma_{\iso}^{\tp}(\Omega_t),
\end{equation}
where $\Omega_t= \set{ w_p \leq t}$.
\end{Lemma}

\begin{proof} Assume that the right-hand side of~\eqref{eq:mass_contro_pcap} is finite, otherwise there is nothing to prove. The function $t\mapsto \abs{\Omega_t\smallsetminus \Crit w_p}$ is monotone continuous in $[0, +\infty)$, hence it is absolutely continuous. The generalised de L'H\^opital rule gives
\begin{align}
 \liminf_{t \to +\infty} \ma_{\iso}^{\tp}(\Omega_t)&{}\geq\liminf_{t \to +\infty} \frac{1}{2p\pi \ncapa_p(\partial \Omega_t)^{\frac{2}{3-p}}}\bigg( \abs{\Omega_t\smallsetminus \Crit w_p } -\frac{4 \pi}{3} \ncapa_p (\partial \Omega_t)^\frac{3}{3-p}\bigg)\nonumber\\
 &{}\geq \liminf_{t \to +\infty} \frac{(3-p)}{4 p\pi \ncapa_p(\partial \Omega_t)^{\frac{2}{3-p}}}\bigg(\int_{\partial \Omega_t} \frac{1}{\abs{\D w_p}} \dif \sigma-\frac{4 \pi}{3-p} \ncapa_p (\partial \Omega_t)^\frac{3}{3-p}\bigg).\label{eq:potential-delhopital}
\end{align}
By H\"{o}lder inequality, we have that
\[
 \int_{\partial \Omega_t }\frac{1}{\abs{\D w_p}} \dif \sigma \geq \bigg(\int_{\partial \Omega_t}\abs{\D w_p}^2 \dif \sigma\bigg)^{-\frac{p}{3-p}}\big( 4 \pi (3-p)^{p-1}\ncapa_p(\partial \Omega_t)\big)^{\frac{3}{3-p}}.
\]
Plugging it in~\eqref{eq:potential-delhopital}, we obtain
\[
 \liminf_{t \to +\infty} \ma_{\iso}^{\tp}(\Omega_t)\geq\liminf_{t \to +\infty} \frac{\ncapa_p(\partial \Omega_t)^{\frac{1}{3-p}}}{p\Big(\int_{\partial \Omega_t} \frac{\abs{ \D w_p}^2}{(3-p)^2} \dif \sigma\Big)^{\frac{p}{3-p}} }\left[ ( 4\pi)^{\frac{p}{3-p}}-\Bigg(\int_{\partial \Omega_t}\frac{\abs{ \D w_p}^2}{(3-p)^2} \dif \sigma \Bigg)^{\frac{p}{3-p}}\right].
\]
To simplify the notation, denote $f(z)=z^{p/(3-p)}$ and $z(t)=\int_{\partial \Omega_t} \abs{\D w_p}^2 /(3-p)^2 \dif \sigma$. Since $z(t)\leq 4 \pi$ in virtue of our assumptions, by Lemma~\ref{thm:inequality_modifiedgeroch_AMMO}
\begin{align}
 \liminf_{t \to +\infty} \ma_{\iso}^{\tp}(\Omega_t)&{}\geq \liminf_{t \to +\infty}\frac{1}{p f(z(t))} \frac{f(4 \pi )-f(z(t))}{4\pi-z(t)} \ncapa_p(\partial \Omega_t)^{\frac{1}{3-p}}( 4 \pi -z(t)) \nonumber\\
 &{}=\liminf_{t \to +\infty}\frac{4\pi(3-p)}{p f(z(t))} \frac{f(4 \pi )-f(z(t))}{4\pi-z(t)} \tilde{\ma}_{H}^{\tp}(\partial \Omega_t) \nonumber\\
 &{}\geq \liminf_{t \to +\infty}\frac{4\pi(3-p)}{p f(z(t))} \frac{f(4 \pi )-f(z(t))}{4\pi-z(t)} \ma_{H}^{\tp}(\partial \Omega_t).\label{eq:function_NPT_nonnegative}
\end{align}
The theorem follows once we prove the following claim.
\begin{Claim}
There exists a divergent increasing sequence $(t_n)_{n \in \N}$ realising the rightmost limit inferior of~\eqref{eq:function_NPT_nonnegative} and such that $z(t_n)\to 4 \pi$ as $n\to +\infty$.
\end{Claim}

Indeed, we would have
\[
 \lim_{n \to +\infty} \frac{4\pi}{f(z(t_n))} =(4 \pi )^{\frac{3-2p}{3-p}},\qquad \lim_{n\to +\infty}\frac{f(4 \pi )-f(z(t_n))}{4\pi-z(t_n)} = f'(4 \pi ) = \frac{p}{3-p} (4\pi)^{\frac{2p-3}{3-p}},
\]
that plugged into~\eqref{eq:function_NPT_nonnegative}, gives~\eqref{eq:mass_contro_pcap} in virtue of Theorem~\ref{thm:AMMO_monotonicity} and Lemma~\ref{thm:inequality_modifiedgeroch_AMMO}.

Let $t_n$ be a divergent increasing sequence $(t_n)_{n \in \N}$ realising the rightmost limit inferior of~\eqref{eq:function_NPT_nonnegative}. By Lemma~\ref{thm:inequality_modifiedgeroch_AMMO}, we have two possible cases:
\begin{enumerate}\itemsep=0pt
 \item[(1)] there exists $T>0$ such that $\tilde{\ma}^{\tp}_{H}(\partial \Omega_{t_n})\geq 0$ for all $t_n \geq T$, or
 \item[(2)] $\tilde{\ma}^{\tp}_{H}(\partial \Omega_{t_n})< 0$ for all $n\in \N$.
\end{enumerate}

\emph{Case~$1$.} Since $\tilde{\ma}^{\tp}_H(\partial \Omega_{t_n}) \geq 0 $, $z(t_n)\leq 4 \pi $ for every $t_n \geq T$. By contradiction, suppose there exists $\varepsilon>0$ such that $z(t_n) \leq 4 \pi - \varepsilon$ for every $n$ sufficiently large. Then, by~\eqref{eq:function_NPT_nonnegative}, there exists $\kst(p,\varepsilon)>0$ such that
\[
 +\infty>\liminf_{t\to +\infty} \ma^{\tp}_{\iso}(\Omega_{t}) \geq \lim_{n \to +\infty} \kst(p,\varepsilon)\ncapa_{p}(\partial \Omega_{t_n})^{\frac{1}{3-p}},
\]
which is clearly a contradiction. Hence, up to a not relabeled subsequence, $z(t_n) \to 4\pi$ as $n \to +\infty$. This proves the claim in this case.

\emph{Case~$2$.} %\label{case:2}
Since $\tilde{\ma}^{\tp}_{H}(\partial \Omega_{t_n}) <0 $, $z(t_n) \geq 4 \pi$ for every $n\in \N$. Suppose by contradiction $z(t_n) \geq 4 \pi + \varepsilon$ for some $\varepsilon>0$. Then, by Theorem~\ref{thm:AMMO_monotonicity}, there exists $\kst(p,\varepsilon)>0$ such that
\[
 \tilde{\ma}^{(p)}_{H}(\partial \Omega) \leq \lim_{t \to +\infty} \tilde{\ma}^{\tp}_{H}(\partial \Omega_t) \leq -\kst(p,\varepsilon)\lim_{n \to +\infty}\ncapa_{p}(\partial \Omega_{t_n})^{\frac{1}{3-p}}=-\infty,
\]
which is a contradiction since $\abs{\D w_p} \in \CS^0(\partial \Omega)$, proving the claim also in this case.\end{proof}

\begin{proof}[Author proof, nonempty-boundary case]\setcounter{equation}{7}
% Original source lines 659--667
 We do now treat the case $\partial M \neq \varnothing$. Let $w_p$ the solution to~\eqref{eq:pb-intro} and define $\Omega_t = \set{w_p \leq t}$. Observe that we can write
 \[
 2\ncapa_p(\partial M)^{\frac{1}{(3-p)}} = 2 \ma_{H}^{\tp}(\partial M)+(3-p)\tilde{\ma}_{H}^{\tp}(\partial M).
 \]
 Then, combining the monotonicity of $\ma_{H}^{\tp}$, that of $\tilde{\ma}_{H}^{\tp}$ following from~\eqref{eq:weak_derivative}, and the asymptotic comparison Lemma~\ref{thm:mass_control_pcap}, we obtain
\begin{align}
 2\ncapa_p(\partial M)^{\frac{1}{(3-p)}} &{}\leq \lim_{t \to +\infty} 2 \ma_{H}^{\tp}(\partial \Omega_t)+(3-p)\tilde{\ma}_{H}^{\tp}(\partial \Omega_t)\nonumber\\
 &{} \leq \liminf_{t \to +\infty} (5-p) \ma_{\iso}^{\tp}(\Omega_t) \leq (5-p) \ma^{\tp}_{\iso}.\label{eq:monotonicitiescombined}
\end{align}

\end{proof}
% Original source lines 931--1090
\appendix

\section[Monotonicities along the p-capacitary potential]{Monotonicities along the $\boldsymbol{p}$-capacitary potential}
\label{app:monotonicities}
Here we slightly improve the monotonicity results in~\cite{agostiniani_riemannianpenroseinequalitynonlinear_2022,chan_monotonicitygreenfunctions_2022}. Inspired by these two works, we are approximating the $p$-capacitary potential with a family of smooth functions. To enter more in detail, let $(M,g)$ be a strongly $p$-nonparabolic Riemannian manifold with (possibly empty) boundary. Let $\Omega \subset M$ be homologous to $\partial M$ and $u_p$ is the solution to~\eqref{eq:pb-intro} starting at $\Omega$. For every $T>1$ let $\Omega_T$ be strictly homologous to $\partial M$ with connected boundary and containing $\set{u_p>\alpha_p(T)}$, where $\alpha_p(T) = T^{-(3-p)/(p-1)}$. Then, we define $u_p^\varepsilon$ as the solution to the following boundary value problem:
\begin{equation}\label{eq:pb-intro-eps}
 \begin{cases}
 \Delta_{p}^\varepsilon u^\varepsilon_p = 0 &\text{on $\Int\Omega_T \smallsetminus \Omega $,}\\
 u_p^\varepsilon=1 & \text{on $\partial \Omega$,}\\
 u_p^\varepsilon=u_p & \text{on $\partial\Omega_T$,}
 \end{cases}
\end{equation}
where
\[
 \Delta_{p}^\varepsilon f = \div\big( \abs{ \D f}_\varepsilon^{p-2} \D f\big) \qquad \text{and} \qquad {\abs{}}_\varepsilon= \sqrt{ \abs{}^2 +\varepsilon^2}.
\]
The function $u^\varepsilon_p$ is smooth away from the exterior boundary and converges in $\CS^{1,\beta}_{\loc}$ to the $p$-capacitary potential $u_p$ as $\varepsilon \to 0^+$. Indeed, this family was used in~\cite{dibenedetto_alphalocalregularityweak_1983} to prove $\CS^{1,\beta}_{\loc}$-regularity of $p$-harmonic functions. Moreover, looking more carefully at the proof of~\cite[Lemma~2.1]{lou_singularsetslocalsolutions_2008}, $\abs{\D u^\varepsilon_p}^{p-1}$~is uniformly bounded in $W^{1,2}_{\loc}$. Hence, up to a not relabeled subsequence, we can always assume that $\abs{ \D u_p^\varepsilon}^{p-1}$ weakly converges in $W^{1,2}_{\loc}$. Moreover, since $\abs{ \D u_p^\varepsilon}$ converges uniformly to~$| \D u_p |$, the weak limit of $\D \abs{ \D u_p^\varepsilon}^{p-1}$ must be $\D | \D u_p |^{p-1}$.

We are now ready to prove Theorem~\ref{thm:AMMO_monotonicity}.

\begin{proof}[Proof of Theorem~\ref{thm:AMMO_monotonicity}] For ease of computations, we rewrite the function $t\mapsto \ma_{H}^{\tp}(\partial \Omega_t)$ in terms of the $p$-capacitary potential, that is,
\[
 U_p(t)= 4 \pi t + \frac{(p-1)^2}{(3-p)^2} t^{\frac{5-p}{p-1}} \int_{\set{u_p = \alpha_p(t)}} | \D u_p |^2 \dif \sigma -\frac{(p-1)}{(3-p)} t^{\frac{2}{p-1}} \int_{\set{u_p = \alpha_p(t)}} | \D u_p |\H \dif \sigma.
\]
Observe that
\[
 \frac{\ncapa_p(\partial \Omega)^\frac{1}{3-p}}{8\pi} U_p\big(\ee^{\frac{t}{3-p}}\big) = \ma_{H}^{\tp}(\partial \Omega_t)= \frac{1}{8\pi}F_p\big(\ncapa_p(\partial \Omega)^\frac{1}{3-p} \ee^{\frac{t}{3-p}}\big),
\]
where $F_p$ is the function defined in~\cite{agostiniani_riemannianpenroseinequalitynonlinear_2022}. Equation~\eqref{eq:AMMO_monotonicity} now follows from computations in~\cite[Section 1.2]{agostiniani_riemannianpenroseinequalitynonlinear_2022}. The function $U_p$ is well defined on $[0,+\infty)$. Indeed, one can observe that
\[
 \abs{ \H} | \D u_p |^2 = | \D u_p |^{3-p} | \D | \D u_p |^{p-1}| \in L^2_{\loc}(M\smallsetminus \Omega),
\]
since $1<p<3$ and $| \D u_p |^{p-1} \in L^\infty_{\loc}\cap W^{1,2}_{\loc} (M\smallsetminus \Omega)$ by~\cite[Lemma~2.1]{lou_singularsetslocalsolutions_2008}. Then, by coarea formula, the function
\[
 t \mapsto \int_{\set{u_p= \alpha_p(t)}} | \D u_p | \H \dif \sigma\in L^1_{\loc}(0,+\infty)
\]
and its equivalence class does not depend on the representative of $| \D u_p | \H $. In particular, the function $t\mapsto \ma_{H}^{\tp}(\partial \Omega_t) \in L^1_{\loc}(0,+\infty)$.

It only remains to prove that it has nonnegative first derivative in the sense of distributions, which both gives that $t\mapsto \ma_{H}^{\tp}(\partial \Omega_t) \in \BV_{\loc}(0,+\infty)$ and that admits a nondecreasing representative. Fix $T>0$ and $u^\varepsilon_p$ be the solution to~\eqref{eq:pb-intro-eps}. One can now define the function
\[
 F_p^\varepsilon(t)= 4 \pi t + \frac{(p-1)^2}{(3-p)^2} t^{\frac{5-p}{p-1}} \int_{\set{u^\varepsilon_p = \alpha_p(t)}} \abs{ \D u^\varepsilon_p}^2 \dif \sigma -\frac{(p-1)}{(3-p)} t^{\frac{2}{p-1}} \int_{\set{u^\varepsilon_p = \alpha_p(t)}} \abs{ \D u^\varepsilon_p}\H^\varepsilon\dif \sigma
\]
for every $t \in (1,T)$, where $\H^\varepsilon$ is the mean curvature of the level $\set{u^\varepsilon_p = \alpha_p(t)}$. By~\cite[Lemma~1.2]{agostiniani_riemannianpenroseinequalitynonlinear_2022}, we have that the function $F^\varepsilon_p$ is almost monotone, in the sense that
\[
 F^\varepsilon_p(t)- F^\varepsilon_p(s) \geq -\varepsilon
 \frac{(p+1)^2(p-1)}{(3-p)^3} \int_{ \set{\alpha_p(t) < u^\varepsilon_p<\alpha_p(s)}}\abs{ \D u_p^\varepsilon}^2(u^\varepsilon_p)^{-\frac{p-1}{3-p}-3} \dif \mu
\]
holds for almost every $0\leq s\leq t \leq T$.
In particular, by coarea formula and Lebesgue differentiation theorem,
\[
 (F^\varepsilon_p)'(t) \geq
 - \varepsilon \frac{(p+1)^2}{(3-p)^2}t^{\frac{2(3-p)}{p-1}} \int_{\set{u^\varepsilon_p=\alpha_p(t)}} \abs{ \D u^\varepsilon_p} \dif \sigma
\]
holds for almost every $t \in (1,T)$.

The monotonicity follows if one can prove the following claim.
\begin{Claim}
$F_p^\varepsilon$ converges to $F_p$ in the sense of distributions as $\varepsilon\to 0^+$.
\end{Claim}
Indeed, if that is the case, for every nonnegative test function $\varphi \in \CS^\infty_c(1,T)$, we obtain that
\begin{align*}
 -\int_1^T \varphi'(t) F_p(t)\dif t & =-\lim_{\varepsilon\to 0} \int_1^T\varphi'(t) F^p_\varepsilon(t)\dif t \\
 & \geq - \frac{(p+1)^2}{(3-p)^2} \lim_{\varepsilon \to 0} \varepsilon \int_1^T \varphi(t) t^{\frac{2(3-p)}{p-1}} \int_{\set{u^\varepsilon_p=\alpha_p(t)}} \abs{ \D u^\varepsilon_p} \dif \sigma\dif t \\
 & =-\frac{(p+1)^2}{(3-p)^2} \lim_{\varepsilon \to 0} \varepsilon \int_{\Omega_T \smallsetminus \Omega} \varphi \big(\alpha_p^{-1}\big(u^\varepsilon_p\big)\big) \frac{ \big| \D u^\varepsilon_p \big|^2}{\big(u^\varepsilon_p\big)^{\frac{3-p}{p-1}+3}} \dif \mu=0,
\end{align*}
since $u_p^\varepsilon$ converges to $u_p$ in $\CS^{1,\beta}_{\loc}$. This shows that $F_p$ has nonnegative first derivative in the sense of distributions, proving its monotonicity.

We now turn to prove the claim. Consider any $\varphi \in \CS^\infty_c(0,+\infty)$. The first term is independent of $\varepsilon$. As far as the second term is concerned, by coarea formula, we have that
\[
 \int_1^T\varphi(t)t^{\frac{5-p}{p-1}} \int_{\set{u_p^\varepsilon=\alpha_p(t)}} \abs{\D u^\varepsilon_p}^2 \dif \sigma \dif t = \frac{(p-1)}{(3-p)} \int_{\Omega_T \smallsetminus \Omega} \big(u^\varepsilon_p\big)^{-\frac{7-p}{3-p}}\varphi \big(\alpha_p^{-1}\big(u^\varepsilon_p\big)\big) \abs{ \D u_p^\varepsilon}^3 \dif \mu.
\]
Since the function $\varphi$ is smooth with compact support in $(1,T)$ and $u_p^\varepsilon$ converges to $u_p$ in $\CS^{1,\beta}_{\loc}$, the right-hand side converges to
\[
 \frac{(p-1)}{(3-p)} \int_{\Omega_T \smallsetminus \Omega} u_p^{-\frac{7-p}{3-p}}\varphi \big(\alpha_p^{-1}\big(u^\varepsilon_p\big)\big)| \D u_p |^3 \dif \mu = \int_1^T \varphi(t)t^{\frac{5-p}{p-1}} \int_{\set{u_p^\varepsilon=\alpha_p(t)}} |\D u_p |^2 \dif \sigma \dif t,
\]
where the identity follows by coarea formula. The last term is a little trickier since it involves second derivatives of the function $u_p^\varepsilon$ that are not converging uniformly as $\varepsilon \to 0$ to the corresponding ones for $u_p$. Employing again the coarea formula and by straightforward computations, we then have that
\begin{gather*}
 \int_1^T\varphi(t)t^{\frac{2}{p-1}} \int_{\set{u_p^\varepsilon=\alpha_p(t)}} \abs{\D u^\varepsilon_p} \H_\varepsilon\dif \sigma \dif t \\
 \qquad {}=
 \int_1^T\varphi(t)t^{\frac{2}{p-1}} \int_{\set{u_p^\varepsilon=\alpha_p(t)}} \frac{\big\langle \D \big| \D u^\varepsilon_p\big|^{p-1}\big|\D u^\varepsilon_p\big\rangle}{\big| \D u_p^\varepsilon\big|^{p-1} }\Bigg( 1 - \frac{(p-2)}{(p-1)} \frac{ \varepsilon^2}{\big|\D u^\varepsilon_p\big|_\varepsilon^2}\Bigg)\dif \sigma \dif t\\
 \qquad {}=\frac{(p-1)}{(3-p)} \int_{\Omega_T \smallsetminus \Omega} \big(u^\varepsilon_p\big)^{-\frac{4}{3-p}}\varphi \big(\alpha_p^{-1}\big(u^\varepsilon_p\big)\big)\frac{\big\langle\D \big| \D u^\varepsilon_p\big|^{p-1}\big|\D u^\varepsilon_p\big\rangle}{\big| \D u_p^\varepsilon\big|^{p-2} }\Bigg( 1 - \frac{(p-2)}{(p-1)} \frac{ \varepsilon^2}{\big|\D u^\varepsilon_p\big|_\varepsilon^2}\Bigg) \dif \mu.
\end{gather*}
As before we want to prove that the right-hand side converges to the corresponding term for $u_p$. Since $u^\varepsilon_p\to u_p$ in $\CS^{1, \beta}_{\loc}$ and $\D\big| \D u^\varepsilon_p\big|^{p-1}\rightharpoonup \D| \D u_p |^{p-1} $ weakly in $L^2_{\loc}$, we have that
\begin{gather*}
 \lim_{\varepsilon \to 0}\frac{(p-1)}{(3-p)} \int_{\Omega_T \smallsetminus \Omega} \big(u^\varepsilon_p\big)^{-\frac{4}{3-p}}\varphi \big(\alpha_p\big(u^\varepsilon_p\big)\big) \frac{\big\langle\D \big| \D u^\varepsilon_p\big|^{p-1}\big|\D u^\varepsilon_p\big\rangle}{\big| \D u_p^\varepsilon\big|^{p-2} }\dif \mu\\
\qquad {} =\frac{(p-1)}{(3-p)} \int_{\Omega_T \smallsetminus \Omega} (u_p)^{-\frac{4}{3-p}}\varphi \big(\alpha_p^{-1}\big(u^\varepsilon_p\big)\big)\frac{\big\langle\D | \D u_p |^{p-1}|\D u_p\big\rangle}{| \D u_p |^{p-2} }\dif \mu\\
\qquad {} =\int_1^T\varphi(t)t^{\frac{2}{p-1}}\int_{\set{u_p=\alpha_p(t)}} |\D u_p | \H \dif \sigma \dif t.
\end{gather*}
Moreover, the remaining term vanishes. H\"older's inequality and equi-boundedness in $L_{\loc}^2(\Omega_T \smallsetminus \Omega)$ of $\big|\D \big|\D u^\varepsilon_p\big|^{p-1}\big|$ yield
\begin{equation}\label{eq:main_inequality_monotonicity_AMMO}
 \int_{K}\abs{\D u^\varepsilon_p}^{3-p}\big|\D \abs{ \D u^\varepsilon_p}^{p-1}\big|\frac{ \varepsilon^2}{\big| \D u^\varepsilon_p\big|_\varepsilon^2}\dif \mu\leq \kst_1\Bigg(\int_K\frac{\varepsilon^4}{\big| \D u^\varepsilon_p\big|_\varepsilon^4} \abs{\D u^\varepsilon_p}^{6-2p}\dif \mu\Bigg)^\frac{1}{2}
\end{equation}
for every $K$ compactly contained in $\set{ \alpha_p(T) < u_p^\varepsilon <1}$ and for some positive constant $\kst_1$. Observe that
\[
 \abs{\D u^\varepsilon_p}^{6-2p}\frac{\varepsilon^4}{\big| \D u^\varepsilon_p\big|_\varepsilon^4} \leq \abs{\D u^\varepsilon_p}^{6-2p}\leq \kst_2,
\]
since the function $\big| \D u_\varepsilon^p\big|$ converges locally uniformly and $1<p<3$. The left-hand side converges almost everywhere to $0$. Indeed, if a point belongs to critical set of $u_p$, \smash{$\big|\D u_\varepsilon^p\big|^{6-2p}\to 0$} as $\varepsilon\to 0$. Otherwise, $\big| \D u_\varepsilon^p\big|$ is definitely bounded away from $0$, then \smash{$\big| \D u^\varepsilon_p\big|_\varepsilon^4$} is not vanishing, thus the left-hand side is controlled by $\varepsilon^4$ up to a constant. By dominated convergence theorem, the right-hand side in~\eqref{eq:main_inequality_monotonicity_AMMO} approaches $0$ as $\varepsilon \to 0$, so does the left-hand side, concluding the step. \end{proof}

We use this theorem to study the quantity introduced in~\cite[Theorem 1.2]{chan_monotonicitygreenfunctions_2022} along the level set of the solution $w_p$ to~\eqref{eq:pb-intro}.

\begin{Theorem}\label{thm:CCLT_monotonicity}
 Let $(M,g)$ be a strongly $p$-nonparabolic Riemannian $3$-manifold with smooth, compact and connected possibly empty boundary $\partial M$. Assume that $H_2(M, \partial M;\Z) = \set{0}$. Let $\Omega\subseteq M$ be bounded closed with connected $\CS^1$-boundary homologous to $\partial M$ and with $\h\in L^2(\partial \Omega)$. Let $w_p$ be the solution to~\eqref{eq:pb-intro} starting at $\Omega$. Then, denoting $\Omega_t = \set{w_p \leq t}$, the function
 \begin{equation}\label{eq:monotoncity_chan}
 t \mapsto \frac{\ncapa_p(\partial \Omega_t)^{-\frac{1}{p-1}}}{4\pi(3-p)}\Bigg( 4 \pi - \int_{\partial \Omega_t} \frac{\abs{\D w_p}^2}{(3-p)^2}\dif \mu\Bigg)
 \end{equation}
 belongs to $W^{1,1}_{\loc}(0,+\infty)$ and
 \begin{equation}\label{eq:weak_derivative}
 \frac{\dif}{\dif t}\Bigg[\ncapa_p(\partial \Omega_t)^{-\frac{1}{p-1}}\Bigg( 4 \pi - \int_{\partial \Omega_t} \frac{\abs{\D w_p}^2}{(3-p)^2}\dif \sigma\Bigg)\Bigg] =- \frac{8\pi}{p-1}\ncapa_p(\partial \Omega_t)^{-\frac{2}{(3-p)(p-1)}}\ma_{H}^{\tp}(\partial \Omega_t),
 \end{equation}
 for almost every $t \in [0,+\infty)$.
\end{Theorem}

\begin{Remark}\label{rmk:cclt_mon}
 Observe that~\eqref{eq:monotoncity_chan} is, up to multiplying by a constant and changing variables, the quantity studied in~\cite{chan_monotonicitygreenfunctions_2022,munteanu_comparisontheorems3dmanifolds_2023} for $p\in(1,2]$ along the level set of the $p$-Green function. In particular, coupled with Theorem~\ref{thm:AMMO_monotonicity}, the above result implies that the monotonicity property of~\eqref{eq:monotoncity_chan} is preserved if, in place of the $p$-Green's function, the $p$-capacitary potential of a connected $\partial \Omega$ with nonnegative $p$-Hawking mass is considered. Clearly, the monotonicity results in~\cite{chan_monotonicitygreenfunctions_2022, munteanu_comparisontheorems3dmanifolds_2023} are recovered applying Theorem~\ref{thm:CCLT_monotonicity} to the $p$-Green's function. Hence, we settle the question raised in~\cite{chan_monotonicitygreenfunctions_2022} about the monotonicity of their quantities in the range $p \in (2, 3)$.

 Finally, observe that $\tilde{\ma}^{\tp}_H$ is obtained multiplying the function~\eqref{eq:monotoncity_chan} by the $p$-capacity term $\ncapa_p(\partial \Omega_t)^{2/(3-p)(p-1)}$ which is exponentially growing as $t \to +\infty$. This term forces the quantity~$\tilde{\ma}^{\tp}_H$ to be monotone \emph{nondecreasing} under assumption~\eqref{eq:energy_hp} (see Lemma~\ref{thm:inequality_modifiedgeroch_AMMO}) even when~\eqref{eq:monotoncity_chan} is monotone \emph{nonincreasing}.
\end{Remark}

\begin{proof}[Proof of Theorem~\ref{thm:CCLT_monotonicity}] Fix $T>1$ and let $u^\varepsilon_p$ the solution to the problem~\eqref{eq:pb-intro-eps}. Consider any $\varphi \in \CS^\infty_c(0,T)$. Employing coarea formula and integration by parts, we have that
\begin{align}
 \int_1^{T} \varphi'(t) \int_{\set{u^\varepsilon_p=\alpha_p(t)}} \abs{ \D u^\varepsilon_p}^2 \dif \sigma \dif t &{}= \frac{(p-1)}{(3-p)} \int_{\Omega_T \smallsetminus \Omega} \varphi' \big(\big(u_p^\varepsilon\big)^{-\frac{p-1}{3-p}} \big)\big(u_p^\varepsilon\big)^{-\frac{2}{3-p}} \abs{ \D u^\varepsilon_p}^3 \dif \mu \nonumber\\
 &{}=- \int_{\Omega_T \smallsetminus \Omega} \big\langle\big| \D u^\varepsilon_p\big|\D u^\varepsilon_p\big| \D \big[\varphi \big(\big(u_p^\varepsilon\big)^{-\frac{p-1}{3-p}}\big)\big]\big\rangle \dif \mu \nonumber\\
 &{}= \int_{\Omega_T \smallsetminus \Omega} \div \big( \abs{ \D u^\varepsilon_p} \D u^\varepsilon_p\big) \varphi \big(\big(u_p^\varepsilon\big)^{-\frac{p-1}{3-p}}\big) \dif \mu. \label{eq:main_chan}
\end{align}
Clearly, employing $\CS^{1, \beta}_{\loc}$ convergence of $u^\varepsilon_p \to u_p$ as $\varepsilon \to 0$, we obtain that
\[
 \lim_{\varepsilon\to 0}\int_1^{T} \varphi'(t) \int_{\set{u^\varepsilon_p=\alpha_p(t)}} \abs{ \D u^\varepsilon_p}^2 \dif \sigma \dif t = \int_1^{T} \varphi'(t) \int_{\set{u_p=\alpha_p(t)}} | \D u_p |^2 \dif \sigma \dif t.
\]
Moreover, a straightforward computation leads to
\[
 \div \big( \abs{ \D u^\varepsilon_p} \D u^\varepsilon_p\big) =\frac{\big| \D u^\varepsilon_p\big|^{2-p}}{(p-1)} \Bigg( (3-p) \frac{ \big| \D u^\varepsilon_p\big|^2}{\big| \D u^\varepsilon_p \big|^2_\varepsilon}- \frac{ \varepsilon^2}{\big| \D u^\varepsilon_p\big|_{\varepsilon}^2} \Bigg) \big\langle\D \big| \D u^\varepsilon_p \big|^{p-1} \big| \D u^\varepsilon_p\big\rangle.
\]
Arguing as in the previous theorem, since $\big|\D u^\varepsilon_p\big| \to |\D u_p |$ locally uniformly and $\D\big|\D u^\varepsilon_p\big|^{p-1} \rightharpoonup\D |\D u_p |^{p-1}$ weakly $L^2_{\loc}$ as $\varepsilon\to 0$, one gets
\begin{gather*}
 \lim_{\varepsilon \to 0} \int_{\Omega_T \smallsetminus \Omega} \frac{\big| \D u^\varepsilon_p\big|^{2-p}}{(p-1)} \Bigg( (3-p) \frac{ \big| \D u^\varepsilon_p\big|^2}{\big| \D u^\varepsilon_p \big|^2_\varepsilon}- \frac{ \varepsilon^2}{\big| \D u^\varepsilon_p\big|_{\varepsilon}^2} \Bigg) \big\langle\D \big| \D u^\varepsilon_p \big|^{p-1} \big| \D u^\varepsilon_p\big\rangle \varphi \big(\big(u_p^\varepsilon\big)^{-\frac{p-1}{3-p}}\big) \dif \mu\\
 \qquad {}=
 \frac{(3-p)}{(p-1)} \int_{\Omega_T \smallsetminus \Omega} \frac{ \big\langle \D | \D u_p |^{p-1} | \D u_p\big\rangle}{| \D u_p |^{p-2}} \varphi \big((u_p)^{-\frac{p-1}{3-p}}\big)\dif \mu \\
 \qquad {}=
 \frac{(3-p)^2}{(p-1)^2}\int_1^{T}\varphi(t) t^{-\frac{2}{p-1}} \int_{\Omega_T \smallsetminus \Omega} \H \abs{ \D u} \dif \sigma \dif t,
\end{gather*}
where the last equality follows by coarea formula and $\H= \big\langle\D | \D u_p |^{p-1} | \D u_p\big\rangle / | \D u_p |^p$. Then, passing to the limit as $\varepsilon \to 0$ in~\eqref{eq:main_chan}, we obtain
\[
 \int_1^{T} \varphi'(t) \int_{\set{u_p=\alpha_p(t)}} | \D u_p |^2 \dif \sigma \dif t =
 \frac{(3-p)^2}{(p-1)^2}\int_1^{T}\varphi(t) t^{-\frac{2}{p-1} } \int_{\set{u_p= \alpha_p(t)}} \H \abs{ \D u} \dif \sigma \dif t.
\]
By arbitrariness of $T$ and $\varphi$ one has that the function
\[
 t \mapsto H_p(t)=\Bigg(4 \pi t^{-\frac{3-p}{p-1}}-\frac{(p-1)^2}{(3-p)^2}t^{\frac{3-p}{p-1}} \int_{\set{u=\alpha_p(t)}} |\D u_p |^2 \dif \sigma\Bigg)
\]
belongs to $W^{1,1}_{\loc}(1,+\infty)$ and its derivative is given by
\begin{align*}
H'_p(t)={}& -\frac{(3-p)}{(p-1)}t^{-\frac{p+1}{p-1}}\Bigg(4 \pi t+ \frac{(p-1)^2}{(3-p)^2} t^{\frac{5-p}{p-1}} \\
&{}\times \int_{\set{u=\alpha_p(t)}} |\D u_p |^2 \dif \sigma - \frac{(p-1)}{(3-p)} t^{\frac{2}{p-1}}\int_{\set{u_p= \alpha_p(t)}} \H \abs{ \D u} \dif \sigma\Bigg)
\end{align*}
holds for almost every $t \in (1,+\infty)$. Then
\begin{align*}
 \frac{\dif}{\dif t}\Bigg[\ncapa_p(\partial \Omega_t)^{-\frac{1}{p-1}}\Bigg( 4 \pi - \int_{\partial \Omega_t} \frac{\big|\D w_p\big|^2}{(3-p)^2}\dif \mu\Bigg)\Bigg] &{}= \frac{\ncapa_p(\partial \Omega)^{-\frac{1}{p-1}}}{3-p} H'_p\big(\ee^{\frac{t}{3-p}}\big) \ee^{\frac{t}{3-p}}\\
 &{}=- \frac{8\pi}{p-1}\ncapa_p(\partial \Omega_t)^{-\frac{2}{(3-p)(p-1)}}\ma_{H}^{\tp}(\partial \Omega_t),
\end{align*}
concluding the proof.\end{proof}
\begin{thebibliography}{99}
\bibitem[1]{agostiniani_sharpgeometricinequalitiesclosed_2020}
Agostiniani V., Fogagnolo M., Mazzieri L., Sharp geometric inequalities for
 closed hypersurfaces in manifolds with nonnegative {R}icci curvature,
 \href{https://doi.org/10.1007/s00222-020-00985-4}{\textit{Invent. Math.}} \textbf{222} (2020), 1033--1101, \href{https://arxiv.org/abs/1812.05022}{arXiv:1812.05022}.

\bibitem[2]{agostiniani_minkowskiinequalitiesnonlinearpotential_2022}
Agostiniani V., Fogagnolo M., Mazzieri L., Minkowski inequalities via nonlinear
 potential theory, \href{https://doi.org/10.1007/s00205-022-01756-6}{\textit{Arch. Ration. Mech. Anal.}} \textbf{244} (2022),
 51--85, \href{https://arxiv.org/abs/1906.00322}{arXiv:1906.00322}.

\bibitem[3]{agostiniani_riemannianpenroseinequalitynonlinear_2022}
Agostiniani V., Mantegazza C., Mazzieri L., Oronzio F., Riemannian {P}enrose
 inequality via nonlinear potential theory, \href{https://arxiv.org/abs/2205.11642}{arXiv:2205.11642}.

\bibitem[4]{agostiniani_monotonicityformulaspotentialtheory_2020}
Agostiniani V., Mazzieri L., Monotonicity formulas in potential theory,
 \href{https://doi.org/10.1007/s00526-019-1665-2}{\textit{Calc. Var. Partial Differential Equations}} \textbf{59} (2020), 6,
 32~pages, \href{https://arxiv.org/abs/1606.02489}{arXiv:1606.02489}.

\bibitem[5]{agostiniani_greenfunctionproofpositive_2021}
Agostiniani V., Mazzieri L., Oronzio F., A {G}reen's function proof of the
 positive mass theorem, \href{https://arxiv.org/abs/2108.08402}{arXiv:2108.08402}.

\bibitem[8]{benatti_monotonicityformulasnonlinearpotential_2022}
Benatti L., Monotonicity formulas in nonlinear potential theory and their
 geometric applications, Ph.D.~Thesis, {U}niversit\`a degli studi di Trento,
 2022.

\bibitem[10]{benatti_minkowskiinequalitycompleteriemannian_2022}
Benatti L., Fogagnolo M., Mazzieri L., Minkowski inequality on complete
 {R}iemannian manifolds with nonnegative {R}icci curvature,
 \href{https://arxiv.org/abs/2101.06063}{arXiv:2101.06063}.

\bibitem[11]{benatti_isoperimetricriemannianpenroseinequality_2022}
Benatti L., Fogagnolo M., Mazzieri L., On the isoperimetric {R}iemannian
 {P}enrose inequality, \href{https://arxiv.org/abs/2212.10215}{arXiv:2212.10215}.

\bibitem[14]{chan_monotonicitygreenfunctions_2022}
Chan P.-Y., Chu J., Lee M.-C., Tsang T.-Y., Monotonicity of the $p$-{G}reen
 functions, \href{https://arxiv.org/abs/2202.13832}{arXiv:2202.13832}.

\bibitem[17]{dibenedetto_alphalocalregularityweak_1983}
DiBenedetto E., {$C^{1+\alpha}$} local regularity of weak solutions of
 degenerate elliptic equations, \href{https://doi.org/10.1016/0362-546X(83)90061-5}{\textit{Nonlinear Anal.}} \textbf{7} (1983),
 827--850.

\bibitem[18]{evans_newprooflocal1_1982}
Evans L.C., A new proof of local {$C^{1,\alpha }$} regularity for solutions of
 certain degenerate elliptic p.d.e., \href{https://doi.org/10.1016/0022-0396(82)90033-X}{\textit{J.~Differential Equations}}
 \textbf{45} (1982), 356--373.

\bibitem[20]{fogagnolo_minimisinghullscapacityisoperimetric_2022}
Fogagnolo M., Mazzieri L., Minimising hulls, {$p$}-capacity and isoperimetric
 inequality on complete {R}iemannian manifolds, \href{https://doi.org/10.1016/j.jfa.2022.109638}{\textit{J.~Funct. Anal.}}
 \textbf{283} (2022), 109638, 49~pages, \href{https://arxiv.org/abs/2012.09490}{arXiv:2012.09490}.

\bibitem[21]{fogagnolo_geometricaspectscapacitarypotentials_2019}
Fogagnolo M., Mazzieri L., Pinamonti A., Geometric aspects of {$p$}-capacitary
 potentials, \href{https://doi.org/10.1016/j.anihpc.2018.11.005}{\textit{Ann. Inst. H.~Poincar\'e C Anal. Non Lin\'eaire}}
 \textbf{36} (2019), 1151--1179, \href{https://arxiv.org/abs/1803.10679}{arXiv:1803.10679}.

\bibitem[22]{heinonen_asuperharmonicfunctionssupersolutionsdegenerate_1988}
Heinonen J., Kilpel\"ainen T., {$A$}-superharmonic functions and supersolutions
 of degenerate elliptic equations, \href{https://doi.org/10.1007/BF02386110}{\textit{Ark. Mat.}} \textbf{26} (1988),
 87--105.

\bibitem[24]{holopainen_nonlinearpotentialtheoryquasiregular_1990}
Holopainen I., Nonlinear potential theory and quasiregular mappings on
 {R}iemannian manifolds, \textit{Ann. Acad. Sci. Fenn. Ser.~A~I~Math.
 Dissertationes} (1990), 1--45.

\bibitem[25]{holopainen_volumegrowthgreenfunctions_1999}
Holopainen I., Volume growth, {G}reen's functions, and parabolicity of ends,
 \href{https://doi.org/10.1215/S0012-7094-99-09714-4}{\textit{Duke Math.~J.}} \textbf{97} (1999), 319--346.

\bibitem[27]{huisken_inversemeancurvatureflow_2001}
Huisken G., Ilmanen T., The inverse mean curvature flow and the {R}iemannian
 {P}enrose inequality, \href{https://doi.org/10.4310/jdg/1090349447}{\textit{J.~Differential Geom.}} \textbf{59} (2001),
 353--437.

\bibitem[28]{jauregui_admmasscapacityvolumedeficit_2020}
Jauregui J.L., {ADM} mass and the capacity-volume deficit at infinity,
 \href{https://arxiv.org/abs/2002.08941}{arXiv:2002.08941}.

\bibitem[30]{kotschwar_localgradientestimatesharmonic_2009}
Kotschwar B., Ni L., Local gradient estimates of {$p$}-harmonic functions,
 {$1/H$}-flow, and an entropy formula, \href{https://doi.org/10.24033/asens.2089}{\textit{Ann. Sci. \'Ec. Norm.
 Sup\'er.~(4)}} \textbf{42} (2009), 1--36, \href{https://arxiv.org/abs/0711.2291}{arXiv:0711.2291}.

\bibitem[31]{ladyzenskaja_linearquasilinearellipticequations_1968}
Ladyzhenskaya O.A., Ural'tseva N.N., Linear and quasilinear elliptic equations,
 Academic Press, New York, 1968.

\bibitem[32]{lewis_regularityderivativessolutionscertain_1983}
Lewis J.L., Regularity of the derivatives of solutions to certain degenerate
 elliptic equations, \href{https://doi.org/10.1512/iumj.1983.32.32058}{\textit{Indiana Univ. Math.~J.}} \textbf{32} (1983),
 849--858.

\bibitem[33]{lieberman_boundaryregularitysolutionsdegenerate_1988}
Lieberman G.M., Boundary regularity for solutions of degenerate elliptic
 equations, \href{https://doi.org/10.1016/0362-546X(88)90053-3}{\textit{Nonlinear Anal.}} \textbf{12} (1988), 1203--1219.

\bibitem[34]{lou_singularsetslocalsolutions_2008}
Lou H., On singular sets of local solutions to {$p$}-{L}aplace equations,
 \href{https://doi.org/10.1007/s11401-007-0312-y}{\textit{Chinese Ann. Math. Ser.~B}} \textbf{29} (2008), 521--530.

\bibitem[35]{mari_flowlaplaceapproximationnew_2022}
Mari L., Rigoli M., Setti A.G., On the {$1/H$}-flow by {$p$}-{L}aplace
 approximation: new estimates via fake distances under {R}icci lower bounds,
 \href{https://doi.org/10.1353/ajm.2022.0016}{\textit{Amer.~J.~Math.}} \textbf{144} (2022), 779--849, \href{https://arxiv.org/abs/1905.00216}{arXiv:1905.00216}.

\bibitem[36]{moser_inversemeancurvatureflow_2007}
Moser R., The inverse mean curvature flow and {$p$}-harmonic functions,
 \href{https://doi.org/10.4171/JEMS/73}{\textit{J.~Eur. Math. Soc.}} \textbf{9} (2007), 77--83.

\bibitem[37]{moser_inversemeancurvatureflow_2008}
Moser R., The inverse mean curvature flow as an obstacle problem,
 \href{https://doi.org/10.1512/iumj.2008.57.3385}{\textit{Indiana Univ. Math.~J.}} \textbf{57} (2008), 2235--2256.

\bibitem[38]{munteanu_comparisontheorems3dmanifolds_2023}
Munteanu O., Wang J., Comparison theorems for 3{D} manifolds with scalar
 curvature bound, \href{https://doi.org/10.1093/imrn/rnab307}{\textit{Int. Math. Res. Not.}} \textbf{2023} (2023),
 2215--2242, \href{https://arxiv.org/abs/2105.12103}{arXiv:2105.12103}.

\bibitem[39]{ni_meanvaluetheoremsmanifolds_2007}
Ni L., Mean value theorems on manifolds, \href{https://doi.org/10.4310/AJM.2007.v11.n2.a6}{\textit{Asian~J.~Math.}} \textbf{11}
 (2007), 277--304, \href{https://arxiv.org/abs/math.DG/0608608}{arXiv:math.DG/0608608}.

\bibitem[41]{serrin_localbehaviorsolutionsquasilinear_1964}
Serrin J., Local behavior of solutions of quasi-linear equations, \href{https://doi.org/10.1007/BF02391014}{\textit{Acta
 Math.}} \textbf{111} (1964), 247--302.

\bibitem[42]{tolksdorf_dirichletproblemquasilinearequations_1983}
Tolksdorf P., On the {D}irichlet problem for quasilinear equations,
 \href{https://doi.org/10.1080/03605308308820285}{\textit{Comm. Partial Differential Equations}} \textbf{8} (1983), 773--817.

\bibitem[43]{trudinger_harnacktypeinequalitiestheir_1967}
Trudinger N.S., On {H}arnack type inequalities and their application to
 quasilinear elliptic equations, \href{https://doi.org/10.1002/cpa.3160200406}{\textit{Comm. Pure Appl. Math.}} \textbf{20}
 (1967), 721--747.

\bibitem[44]{uraltseva_degeneratequasilinearellipticsystems_1968}
Ural'tseva N.N., Degenerate quasilinear elliptic systems, \textit{Zap. Nauchn.
 Sem. Leningrad. Otdel. Mat. Inst. Steklov. (LOMI)} \textbf{7} (1968),
 184--222.

\bibitem[45]{valtorta_laplaceoperatorriemannianmanifolds_2013}
Valtorta D., On the $p$-{L}aplace operator on {R}iemannian manifolds, Ph.D.~Thesis, {U}niversit\`a degli Studi di Milano, 2013, \href{https://arxiv.org/abs/1212.3422}{arXiv:1212.3422}.

\bibitem[47]{xia_anisotropiccapacityanisotropicminkowski_2022}
Xia C., Yin J., The anisotropic {$p$}-capacity and the anisotropic {M}inkowski
 inequality, \href{https://doi.org/10.1007/s11425-021-1884-1}{\textit{Sci. China Math.}} \textbf{65} (2022), 559--582,
 \href{https://arxiv.org/abs/2012.13933}{arXiv:2012.13933}.

\bibitem[49]{zhu_comparisongeometryriccicurvature_1997}
Zhu S., The comparison geometry of {R}icci curvature, in Comparison {G}eometry
 ({B}erkeley, {CA}, \textit{Math. Sci. Res. Inst. Publ.}, Vol.~30, Cambridge
 University Press, Cambridge, 1997, 221--262.

\end{thebibliography}
\end{document}
